\chapter{Estado del arte}

La aplicación de la inteligencia artificial a la resonancia magnética (RM)
cerebral en neurooncología ha experimentado un crecimiento notable, impulsada por
la disponibilidad de \emph{datasets} públicos, el avance de las arquitecturas de
aprendizaje profundo y la necesidad de herramientas de apoyo clínico útiles y
realistas. Buena parte de la literatura se ha centrado, sin embargo, en tareas
muy específicas ---segmentación exhaustiva del tumor, clasificación histológica o
predicción de biomarcadores moleculares--- que, pese a su valor científico, no
siempre se traducen de forma directa en aplicaciones integrables en el flujo
asistencial cotidiano. La priorización automática de estudios con sospecha de
masa tumoral, en cambio, constituye una tarea de interés asistencial directo y
metodológicamente más acotada.

Esta revisión se articula en torno a seis ejes: el contexto clínico y el
diagnóstico por imagen de las masas tumorales intracraneales; el papel de la RM y
las limitaciones del diagnóstico radiológico; el aprendizaje profundo en imagen
médica como sistema de apoyo; sus enfoques, datasets y retos de generalización en
RM cerebral; el sesgo de dominio y el aprendizaje por atajos como amenaza a la
validez; y, por último, los sistemas de priorización en radiología. El recorrido
converge en una observación que vertebra el resto de la memoria: cuando se
construyen sistemas a partir de datos públicos heterogéneos, la \emph{evaluación
rigurosa del sesgo de dominio} deja de ser un detalle metodológico para
convertirse en el problema central, y es ahí donde se sitúa la contribución de
este trabajo.

\section{Diagnóstico por imagen de masas tumorales intracraneales}

\subsection{Relevancia clínica de las masas tumorales intracraneales}

Las masas tumorales intracraneales constituyen un grupo de lesiones de elevada
relevancia clínica y epidemiológica por su impacto neurológico, diagnóstico y
terapéutico~\cite{ostrom2023cbtrus,mcneill2016epidemiology}. Pueden diferir de
manera significativa en su origen, comportamiento biológico, velocidad de
crecimiento, capacidad infiltrativa y pronóstico, diferencias que se recogen en
los sistemas de clasificación histológica y molecular de referencia~\cite{louis2021who,who2007cns,kleihues2000who}.
Algunas presentan una evolución relativamente lenta, mientras que otras ---de
forma señalada los gliomas de alto grado--- muestran un comportamiento agresivo,
progresivo e infiltrativo~\cite{wen2008gliomas}. En todos los casos, su
localización dentro del cráneo confiere especial relevancia a su detección, ya
que incluso lesiones histológicamente poco agresivas pueden producir síntomas por
efecto masa, edema o compresión de estructuras críticas~\cite{contreras2017epidemiologia,perezsegura2025tumores}.

En la práctica clínica, el hallazgo de una masa intracraneal suele desencadenar
una cadena de decisiones de gran importancia: estudios complementarios,
valoración por neurocirugía, discusión multidisciplinar y eventual inicio de
tratamiento. Por ello, la identificación temprana de estudios con sospecha de
lesión ocupante de espacio puede tener repercusiones directas sobre los tiempos
de respuesta y la organización del circuito asistencial. A los efectos de este
trabajo no es necesario profundizar en una clasificación exhaustiva de subtipos
ni en su estratificación molecular: lo esencial es reconocer que estas masas
constituyen un grupo clínicamente prioritario cuya sospecha en una RM puede
justificar una revisión más temprana, situando el problema en un plano operativo
---priorizar, no sustituir el diagnóstico especializado.

\subsection{Complejidad diagnóstica y necesidad de identificación temprana}

El diagnóstico de una masa intracraneal representa un reto importante por la
heterogeneidad de presentación de estas lesiones y por la complejidad de su
caracterización mediante imagen. La sintomatología es variable e inespecífica
---cefalea, focalidad neurológica, crisis epilépticas, alteraciones cognitivas o
hallazgos incidentales---, lo que confiere a la RM cerebral un papel central en
la detección inicial y en la evaluación posterior~\cite{cha2009neuroimaging}.

La necesidad de identificación temprana no debe entenderse aquí como cribado
poblacional, sino como la revisión temprana de estudios ya indicados clínicamente: el problema
no es decidir si una RM debe realizarse, sino gestionar de forma eficiente el
orden en que los estudios obtenidos son revisados. Cuando una exploración
contiene una masa tumoral, su identificación rápida puede condicionar el
itinerario del paciente y acelerar decisiones de alto impacto. De ahí el interés,
desde una perspectiva organizativa, por sistemas capaces de señalar
automáticamente los estudios con mayor probabilidad de contener una lesión
ocupante de espacio. Este planteamiento es especialmente pertinente en contextos
de elevada carga asistencial ---guardias, turnos nocturnos, servicios con alta
demanda---, y presenta una ventaja conceptual: permite formular una tarea
clínicamente significativa sin exigir que el modelo emita un diagnóstico completo,
acotando la pregunta a si un estudio merece revisión prioritaria por su posible
contenido tumoral.

\section{RM cerebral y limitaciones del diagnóstico radiológico}

\subsection{Papel de la resonancia magnética en neurooncología}

La RM cerebral se ha consolidado como técnica de imagen de referencia en
neurooncología por su elevada resolución espacial y su capacidad para diferenciar
los tejidos blandos del sistema nervioso central~\cite{villanuevameyer2017imaging}.
A diferencia de otras modalidades, permite una evaluación detallada de la anatomía
cerebral y de las alteraciones asociadas a una masa ---edema, necrosis, efecto
masa o alteración de la barrera hematoencefálica--- y se emplea tanto en la
detección inicial como en la estimación de la extensión, la planificación
terapéutica y el seguimiento~\cite{martucci2023mri}. Las distintas secuencias
aportan información complementaria; las secuencias estructurales convencionales
ocupan un lugar central por su disponibilidad y relevancia diagnóstica, y la
caracterización multiparamétrica resulta especialmente útil en determinadas
lesiones~\cite{laukamp2021meningioma}.

Desde el punto de vista del aprendizaje automático, la RM ofrece una fuente de
información rica pero introduce retos metodológicos. En la literatura basada en
datasets públicos es habitual restringir las modalidades a las comunes entre
cohortes para favorecer la comparabilidad: aunque clínicamente sería deseable
trabajar con protocolos multiparamétricos completos, la disponibilidad desigual
de secuencias como T1 con contraste o FLAIR entre cohortes puede introducir
sesgos si su presencia queda asociada a una clase o a un origen concreto. Por
ello, en términos generales, la elección de las modalidades de entrada no es una
cuestión meramente técnica sino una decisión metodológica con impacto sobre la
validez del experimento, y restringirse a modalidades estructurales comunes como
T1 y T2 suele interpretarse como una estrategia conservadora pero sólida, que
sacrifica riqueza clínica a cambio de una comparación más limpia entre dominios.

\subsection{Información visual relevante y dificultad de interpretación}

La RM cerebral contiene abundante información visual útil para sospechar una masa
intracraneal: lesión ocupante de espacio, alteración de la arquitectura
anatómica, edema perilesional, heterogeneidad de señal, necrosis, efecto masa y
desplazamiento de estructuras adyacentes. Su integración permite al radiólogo
formular una impresión diagnóstica inicial. No obstante, la interpretación de
estos patrones no es trivial: algunas lesiones tumorales comparten características
radiológicas con entidades no neoplásicas ---procesos inflamatorios, infecciosos
o cambios postratamiento--- y lesiones de distinta biología pueden mostrar
hallazgos parcialmente solapados, lo que limita la especificidad de la
técnica~\cite{villanuevameyer2017imaging}.

Esta dificultad interpretativa refuerza la pertinencia de un enfoque de
priorización frente a uno de diagnóstico exhaustivo: la RM, aun sin resolver por
sí sola todo el diagnóstico diferencial, contiene información suficiente para
sospechar que ciertos estudios merecen revisión temprana. Esta observación es
relevante al plantear tareas binarias de clasificación a nivel de estudio,
centradas en la presencia o ausencia de masa, que pueden apoyarse en la estructura
visual global de la exploración sin exigir anotaciones a nivel de vóxel.

\subsection{Limitaciones del diagnóstico radiológico y papel del radiólogo}

Pese a su valor, la interpretación radiológica presenta limitaciones. La falta de
especificidad de muchos hallazgos hace que distintas entidades compartan patrones
visuales, y la naturaleza infiltrativa de muchos tumores ---en especial los
gliomas de alto grado--- dificulta delimitar sus márgenes y puede llevar a
infraestimar la extensión real en la imagen convencional~\cite{wen2008gliomas}. A
ello se añade el componente de subjetividad: la experiencia del radiólogo influye
en la lectura y la variabilidad interobservador puede afectar a la consistencia
de los informes, especialmente en estudios complejos. El incremento sostenido del
volumen de exploraciones agrava esta carga.

En este contexto, el radiólogo sigue siendo el profesional central e
insustituible. Por ello, las herramientas de inteligencia artificial de mayor
interés clínico no son las que aspiran a reemplazar su juicio, sino las que lo
apoyan de forma concreta; entre ellas, la priorización automática de estudios
sospechosos destaca como estrategia para redistribuir la carga asistencial y
mejorar los tiempos de respuesta.

\section{Deep Learning en imagen médica}

\subsection{Fundamentos y aplicaciones generales}

La inteligencia artificial es una de las áreas de mayor crecimiento en imagen
médica de la última década, favorecida por la digitalización masiva de los
estudios, el aumento de la capacidad computacional y la disponibilidad de
algoritmos capaces de aprender patrones complejos a partir de grandes volúmenes
de datos~\cite{litjens2017survey,avanzo2024evolution}. En este marco, las redes
neuronales convolucionales se han consolidado como una de las arquitecturas más
relevantes para el análisis de imagen, por su capacidad de extraer
automáticamente características espaciales discriminativas, con utilidad
demostrada en detección de lesiones, clasificación de estudios, segmentación de
estructuras y predicción de resultados clínicos~\cite{litjens2017survey,langs2023ai}.

El interés de estas técnicas no reside únicamente en su capacidad técnica: su
valor asistencial depende de que resuelvan problemas concretos, clínicamente
relevantes y bien planteados~\cite{driver2019radiology}. De ahí el interés
creciente por aplicaciones integrables de forma realista en el flujo de trabajo
hospitalario, más allá de tareas altamente sofisticadas.

\subsection{Sistemas de apoyo al diagnóstico}

El impacto clínico de la inteligencia artificial en radiología depende de su
integración como sistema de apoyo, complementario al radiólogo y no sustitutivo.
Su potencial radica en mejorar la eficiencia, la consistencia y la calidad del
proceso diagnóstico~\cite{sharma2020diagnostic}. Estos sistemas adoptan formas
diversas ---detección asistida por ordenador, clasificadores probabilísticos o
mecanismos de priorización---, con evidencia de adopción en distintos dominios de
la radiología, como la imagen de mama o el diagnóstico de
fracturas~\cite{hickman2021breast,zech2022fracture,yeasmin2024advances}; en
neurorradiología, las aproximaciones radiómicas y de aprendizaje profundo se han
incorporado igualmente al análisis de imagen~\cite{wagner2021radiomics}. La
priorización automática, que se desarrolla en la Sección~\ref{sec:triaje},
constituye una de estas formas de apoyo, especialmente atractiva en entornos de
alta carga asistencial.

\section{Deep Learning aplicado a RM cerebral: enfoques, datasets y retos de generalización}

\subsection{Principales enfoques: segmentación, clasificación y tareas relacionadas}

La literatura sobre aprendizaje profundo en RM cerebral de tumores ha explorado
múltiples tareas. La más desarrollada es la segmentación automática del tumor,
impulsada de forma decisiva por competiciones internacionales y por la
disponibilidad de anotaciones manuales detalladas~\cite{labella2023brats}, que ha
permitido desarrollar arquitecturas 3D capaces de delimitar regiones tumorales
con elevada precisión~\cite{kamnitsas2017deepmedic}. Junto a ella se han
desarrollado tareas de clasificación ---detección de presencia tumoral,
diferenciación de grados o subtipos--- y enfoques de predicción de biomarcadores,
pronóstico o respuesta al tratamiento~\cite{litjens2017survey}.

No todas estas tareas tienen la misma proximidad al flujo asistencial. Mientras
que la segmentación exhaustiva o la predicción molecular exigen niveles elevados
de anotación, validación e integración, las tareas binarias a nivel de estudio
---orientadas a detectar presencia tumoral o a priorizar estudios sospechosos---
pueden constituir un punto de partida más viable, en especial cuando se trabaja
con datasets públicos heterogéneos y recursos computacionales razonables.

\subsection{Datasets públicos y retos de generalización}

El desarrollo de modelos en tumores cerebrales ha dependido estrechamente de
datasets públicos estandarizados, que se han consolidado como referencia del
campo. Esta dependencia ha puesto de manifiesto, sin embargo, una limitación
importante: el rendimiento alcanzado en bases de datos controladas no siempre se
traslada a entornos clínicos reales. Las diferencias en protocolos de
adquisición, resolución espacial, intensidad de campo, preprocesado y
características demográficas introducen variabilidad suficiente como para
comprometer la robustez de los modelos. Por ello, la validación externa y la
evaluación \emph{cross-dataset} adquieren un papel central: no basta con obtener
buenos resultados dentro del dominio de entrenamiento, sino que es necesario
estudiar en qué medida el modelo mantiene capacidad discriminativa al aplicarse a
datos de otra procedencia; desde una perspectiva traslacional, la robustez entre
dominios es una condición indispensable para considerar útil cualquier sistema de
apoyo.

Esta cuestión convierte el problema en un objeto de análisis metodológicamente
rico. La utilidad de un modelo no depende solo de su discriminación global, sino
también del compromiso entre sensibilidad y especificidad, del efecto de los
umbrales de decisión, de la estabilidad del rendimiento entre dominios y de la
calibración de las probabilidades emitidas. Desde una perspectiva de ingeniería
matemática, este conjunto de preguntas hace de la evaluación cuidadosa del modelo
una parte esencial de la aportación, y no un mero apéndice de la implementación.
Ahora bien, cuando la clase positiva y la negativa proceden de datasets distintos,
la generalización deja de ser solo un reto de robustez y se convierte en una
amenaza directa a la validez, que se aborda en la sección siguiente.

\section{Sesgo de dominio, shortcut learning y fugas de información en Deep Learning médico}

\subsection{Aprendizaje por atajos y correlaciones espurias}

Un fenómeno transversal al aprendizaje profundo es el aprendizaje por atajos
(\emph{shortcut learning}): la tendencia de las redes a resolver una tarea
mediante reglas que funcionan en los datos de entrenamiento pero no se
corresponden con la solución pretendida ni generalizan a condiciones
nuevas~\cite{geirhos2020shortcut}. En imagen médica, este fenómeno se manifiesta
con frecuencia como una correlación espuria entre la etiqueta y la procedencia de
los datos. Dos casos se han convertido en referentes: modelos que aparentaban
detectar neumonía en radiografía de tórax y que, en realidad, explotaban
marcadores asociados al hospital o al equipo de adquisición~\cite{zech2018pneumonia},
y modelos de detección de COVID-19 que seleccionaban atajos en lugar de
señal patológica genuina~\cite{degrave2021shortcuts}. En ambos, un rendimiento
interno excelente ocultaba una dependencia de variables ajenas a la enfermedad.

\subsection{Efectos de escáner y cambio de dominio}

La raíz de muchas de estas correlaciones espurias está en los efectos de escáner
y de centro: en datos multi-sitio existen diferencias sistemáticas y medibles
---atribuibles al equipo, al protocolo o al preprocesado--- incluso entre sujetos
comparables, que un modelo puede aprender como si fueran señal~\cite{glocker2019scanner}.
El marco causal del cambio de dominio formaliza por qué estas dependencias se
rompen ---o incluso se invierten--- al cambiar de población: cuando una variable
de procedencia actúa como factor de confusión entre la imagen y la etiqueta, abre
un camino predictivo espurio indistinguible del causal mientras no se intervenga
sobre la composición de los datos~\cite{castro2020causality}. El riesgo es máximo
cuando la etiqueta queda correlacionada con el dominio de origen, situación en la
que el modelo puede aprender diferencias de escáner, protocolo, resolución o
distribución de intensidades en lugar de rasgos asociados a la lesión.

\subsection{Fugas de información y validez experimental}

Además del sesgo de dominio, la literatura identifica otras fugas de información
(\emph{leakage}) que inflan artificialmente las métricas y comprometen la validez.
Dos son especialmente relevantes con datos heterogéneos: la fuga por partición,
que se produce cuando muestras de un mismo sujeto aparecen simultáneamente en
entrenamiento y test, y la presencia de duplicados exactos entre particiones.
Como conceptos generales de buena práctica, su control ---particiones a nivel de
sujeto y comprobación de duplicados--- forma parte del diseño experimental
riguroso, y permite distinguir estas fugas de un confound de dominio propiamente
dicho.

Entre las medidas habituales para mitigar estos riesgos figuran el preprocesado
homogéneo de todas las cohortes, la selección de modalidades comunes, la
estandarización espacial e intensiva, los análisis desglosados por procedencia y
la evaluación externa o \emph{cross-dataset}. Sin embargo, la literatura es menos
explícita en \emph{cómo verificar sistemáticamente} si un rendimiento elevado
refleja señal patológica o dependencia del dominio, sobre todo cuando la
correlación entre clase y procedencia es muy alta. Este es precisamente el hueco
que motiva el presente trabajo: la necesidad de un protocolo de auditoría que, de
forma reproducible, ponga a prueba la validez del rendimiento aparente y aísle el
sesgo de dominio de otras explicaciones. Su desarrollo y aplicación se abordan en
los capítulos de diseño experimental y resultados.

\section{Sistemas de priorización/triaje en radiología}
\label{sec:triaje}

La incorporación de la inteligencia artificial en radiología no es solo una
innovación tecnológica, sino una posible respuesta a un problema organizativo
real: el incremento del volumen de estudios y la creciente complejidad de su
interpretación han aumentado la carga de trabajo del radiólogo. Cualquier
herramienta capaz de redistribuir el tiempo de lectura de forma más eficiente
puede tener un impacto relevante, y los sistemas de priorización o \emph{triaje}
se han propuesto específicamente para reordenar la lista de lectura señalando los
estudios potencialmente urgentes~\cite{titano2018cranial}.

En el caso de la RM cerebral, la revisión de estudios sospechosos requiere un
análisis detallado y la integración de múltiples variables; la presión asistencial
hace que la priorización automática resulte especialmente atractiva. Su objetivo
no es reemplazar la lectura médica, sino señalar los casos que merecerían atención
más temprana. Desde una perspectiva organizativa, una de sus ventajas es que
puede aportar valor incluso sin alcanzar un rendimiento compatible con la
automatización diagnóstica: un sistema suficientemente sensible es útil como
mecanismo de alerta o de reorganización de la cola, siempre que sus limitaciones
estén bien comprendidas y la decisión final siga recayendo en el especialista. La
aceptación de estos sistemas, además, depende no solo de su rendimiento técnico,
sino también de su transparencia y de la percepción de utilidad clínica~\cite{lambert2023acceptance,lee2022implementation}.

Esta función específica tiene implicaciones sobre cómo evaluarlos. Frente a los
sistemas de diagnóstico definitivo, que exigen niveles muy altos de especificidad
y validación, los de priorización deben evaluarse en términos operativos:
sensibilidad a umbrales altos, falsos positivos resultantes, fracción de la lista
que podría reorganizarse y consistencia de las probabilidades con el riesgo
real~\cite{sharma2020diagnostic}. Este desplazamiento, desde la precisión técnica
aislada hacia la utilidad operativa, es clave para valorar de forma realista el
interés de la propuesta.

En síntesis, aunque el aprendizaje profundo aplicado a la RM cerebral ha avanzado
notablemente en segmentación, clasificación y predicción molecular, persisten
limitaciones de generalización, validación externa e integración asistencial. La
priorización de estudios con sospecha de masa tumoral ---ya indicados
clínicamente--- es una formulación coherente con las necesidades de la radiología
y motiva el problema de partida de este trabajo. Ahora bien, la revisión
precedente muestra que, al construir estos sistemas con datos públicos
multi-fuente, el riesgo de que la clase quede confundida con la procedencia es
estructural y puede invalidar cualquier métrica interna. Por ello, el núcleo de
la contribución de esta memoria no es un clasificador de triaje desplegable, sino
la \textbf{evaluación rigurosa del sesgo de dominio}: un protocolo de auditoría
que determina si el rendimiento observado refleja detección de lesión o
reconocimiento de la procedencia. Los capítulos siguientes desarrollan los
materiales, la metodología y ese protocolo.
